\chapter{Implementation and Verification of the Spalart-Allmaras Model in FEniCS}
\label{ch:fenics_implementation}

\section{Introduction to FEniCS and the Finite Element Method}
With the theoretical foundation of the Spalart-Allmaras (SA) model established in Chapter~\ref{ch:turbulence_modeling}, this chapter details its numerical translation into a working computational solver. The numerical simulations and eventual topology optimization in this thesis are built upon FEniCS, an open-source computing platform for solving partial differential equations (PDEs). Unlike commercial computational fluid dynamics (CFD) software such as Ansys, which typically relies on the Finite Volume Method (FVM) and provides a "black box" graphical interface, FEniCS uses the Finite Element Method (FEM) and requires the user to explicitly define the mathematical formulation of the problem \cite{Logg2012}. 

To solve a PDE in FEniCS, the governing equations (the "strong form") must be mathematically transformed into a "weak formulation" (or variational form). FEniCS utilizes the Unified Form Language (UFL) to express these weak forms in Python code that closely mirrors the mathematical notation \cite{Langtangen2016}. This high-level mathematical scripting is particularly advantageous for topology optimization, as it allows for the automated derivation of the adjoint equations required for gradient-based optimization algorithms. 

To ensure complete transparency and reproducibility, the entire custom codebase developed for this thesis is available in the public GitHub repository \texttt{TurbulentTO}. To maintain modularity, the repository is divided into two primary directories. The \texttt{TurbulenceModels} directory contains the standalone Spalart-Allmaras fluid dynamics engine and the verification benchmarks discussed throughout this chapter. Conversely, the \texttt{FluidTO} directory houses the optimization algorithms, penalty formulations, and objective functions that will be detailed in Chapters~\ref{ch:topology_optimization_methodology} and \ref{ch:benchmark_results}.

\subsection{Derivation of a Weak Formulation}
Before diving into the complex turbulence equations, it is instructive to demonstrate how a continuous PDE is translated into a discrete FEniCS formulation. To illustrate this methodology, consider the Poisson equation for pressure, which is a fundamental component of incompressible fluid solvers. The strong form of the Poisson equation is given by:
\begin{equation}
	-\nabla^2 p = f \quad \text{in} \, \Omega
\end{equation}
where $p$ is the pressure field, $f$ is a source term (derived from the velocity divergence), and $\Omega$ is the computational domain. To obtain the weak form, the equation is multiplied by a sufficiently smooth test function $v$ and integrated over the entire domain $\Omega$, following standard Galerkin finite element procedures \cite{Logg2012, Langtangen2016}:
\begin{equation}
	-\int_{\Omega} (\nabla^2 p) v \, dx = \int_{\Omega} f v \, dx
\end{equation}
Applying integration by parts (Green's first identity) to the left-hand side yields:
\begin{equation}
	\int_{\Omega} \nabla p \cdot \nabla v \, dx - \int_{\partial \Omega} (\nabla p \cdot \mathbf{n}) v \, ds = \int_{\Omega} f v \, dx
\end{equation}
where $\partial \Omega$ is the boundary of the domain and $\mathbf{n}$ is the outward-pointing normal vector. This integral formulation fundamentally lowers the continuity requirements of the solution, requiring only first-order derivatives to be square-integrable. Furthermore, Neumann boundary conditions (e.g., a prescribed pressure gradient at an outlet) can be naturally substituted directly into the boundary integral $\int_{\partial \Omega} (\nabla p \cdot \mathbf{n}) v \, ds$. This variational approach forms the foundation of all fluid and turbulence equations implemented in this thesis.

\section{Numerical Solution Strategy for Fluid Dynamics}
Solving the full Reynolds-Averaged Navier-Stokes (RANS) equations is numerically challenging due to the non-linear convective terms, the saddle-point nature of the coupled velocity-pressure system, and the intense cross-coupling with the turbulence model. To achieve a stable and efficient steady-state solution in FEniCS that seamlessly translates into the topology optimization framework, the governing equations are rigorously decoupled.

\subsection{Decoupling the Physics: The Picard Iteration Scheme}
Because the momentum equation relies on the turbulent eddy viscosity ($\nu_t$), and the turbulence transport equation relies on the convective velocity ($\mathbf{u}$), solving the entire fluid-turbulence system simultaneously in a massive monolithic matrix frequently leads to severe ill-conditioning and solver divergence. 

To resolve this, the forward physics are decoupled using a \textbf{Picard iteration scheme} (successive substitution), a highly robust approach for segregated RANS solvers \cite{Ferziger2012, Versteeg2007}. During each Picard step, the flow solver advances using the eddy viscosity from the previous iteration. The newly updated velocity field is then passed to the Spalart-Allmaras solver to calculate a new eddy viscosity field. This alternating sequence forms the outer loop of the solver, iteratively updating the decoupled equations until the fully coupled physical system reaches a steady-state equilibrium.

\subsection{Incremental Pressure Correction Scheme (IPCS)}
Within the Picard loop, the Navier-Stokes equations themselves must be solved. Rather than solving for velocity and pressure simultaneously, this implementation utilizes an Incremental Pressure Correction Scheme (IPCS). This algorithmic splitting is fundamentally based on the fractional-step projection methods originally pioneered by Chorin \cite{Chorin1968} and Temam \cite{Temam1969}. To nurse the non-linear convective momentum terms to a stable steady state, a pseudo-time stepping approach ($\Delta t$) is employed strictly for the flow variables. The IPCS algorithm splits the fluid equations into three sequential steps per pseudo-time step:
\begin{enumerate}
	\item \textbf{Tentative Velocity Step:} The momentum equation is solved using the pressure field from the previous iteration ($p^n$) to compute an intermediate, non-divergence-free velocity field ($\mathbf{u}^*$). Because this is a RANS framework, the diffusion term utilizes the effective kinematic viscosity, $(\nu + \nu_t)$. Treating the convective term explicitly and the diffusive term implicitly, the semi-discrete equation is:
	\begin{equation}
		\frac{\mathbf{u}^* - \mathbf{u}^n}{\Delta t} + (\mathbf{u}^n \cdot \nabla) \mathbf{u}^* - \nabla \cdot \left( (\nu + \nu_t^n) \nabla \mathbf{u}^* \right) = -\nabla p^n
	\end{equation}
	
	\item \textbf{Pressure Correction Step:} A Poisson equation is solved to find the pressure variation that enforces the continuity constraint ($\nabla \cdot \mathbf{u}^{n+1} = 0$). The equation for the new pressure field ($p^{n+1}$) is derived as:
	\begin{equation}
		\nabla^2 p^{n+1} = \nabla^2 p^n + \frac{1}{\Delta t} \nabla \cdot \mathbf{u}^*
	\end{equation}
	
	\item \textbf{Velocity Update Step:} The tentative velocity is strictly corrected using the gradient of the pressure difference to yield a mass-conserving velocity field ($\mathbf{u}^{n+1}$) for the current iteration:
	\begin{equation}
		\mathbf{u}^{n+1} = \mathbf{u}^* - \Delta t \nabla (p^{n+1} - p^n)
	\end{equation}
\end{enumerate}
The simulation is deemed fully converged to a steady state when the $L_2$ norm of the residual difference between successive Picard iterations falls below a strictly defined tolerance limit ($\epsilon_{tol}$) for both the fluid and turbulence fields.

\section{Detailed Formulation of the Spalart-Allmaras Model}
As concluded in Chapter~\ref{ch:turbulence_modeling}, the Spalart-Allmaras (SA) one-equation model was selected for its balance of physical accuracy and numerical stability. Unlike the transient flow equations, the SA transport equation is solved in its purely steady-state form within the Picard loop. The steady transport equation for the working variable $\tilde{\nu}$ is given by:
\begin{equation}
	\mathbf{u} \cdot \nabla \tilde{\nu} = c_{b1} \tilde{S} \tilde{\nu} - c_{w1} f_w \left( \frac{\tilde{\nu}}{y} \right)^2 + \frac{1}{\sigma} \left[ \nabla \cdot \left( (\nu + \tilde{\nu}) \nabla \tilde{\nu} \right) + c_{b2} (\nabla \tilde{\nu})^2 \right]
	\label{eq:sa_steady_transport}
\end{equation}
where $\nu$ is the molecular kinematic viscosity, $\sigma = 2/3$ is the turbulent Prandtl number, $c_{b1}$ and $c_{w1}$ are calibrated closure constants, $\tilde{S}$ is a modified vorticity magnitude, $f_w$ is the wall-damping function, and $y$ is the physical distance to the nearest solid boundary \cite{Spalart1994}.

\subsection{Closure Terms and Calibration Constants}
The Spalart-Allmaras model relies on several interconnected closure functions to accurately capture boundary layer physics. The production term relies on a modified vorticity magnitude, $\tilde{S}$, which prevents the destruction of turbulence in regions where the vorticity goes to zero, while maintaining correct scaling near the wall \cite{Spalart1994}. It is defined as:
\begin{equation}
	\tilde{S} = S + \frac{\tilde{\nu}}{\kappa^2 y^2} f_{v2}
\end{equation}
where $S = \sqrt{2 \Omega_{ij}\Omega_{ij}}$ is the scalar magnitude of the vorticity tensor, $\kappa$ is the von Kármán constant, and $y$ is the distance to the nearest wall. The function $f_{v2}$ is a near-wall modifier:
\begin{equation}
	f_{v2} = 1 - \frac{\chi}{1 + \chi f_{v1}} \quad \text{where} \quad \chi = \frac{\tilde{\nu}}{\nu} \quad \text{and} \quad f_{v1} = \frac{\chi^3}{\chi^3 + c_{v1}^3}
\end{equation}

Similarly, the wall destruction term is modulated by the non-dimensional function $f_w$, which accelerates the decay of turbulent viscosity inside the viscous sublayer:
\begin{equation}
	f_w = g \left( \frac{1 + c_{w3}^6}{g^6 + c_{w3}^6} \right)^{1/6}
\end{equation}
where the intermediate variables $g$ and $r$ are evaluated as:
\begin{equation}
	g = r + c_{w2}(r^6 - r) \quad \text{and} \quad r = \min \left( \frac{\tilde{\nu}}{\tilde{S} \kappa^2 y^2}, 10 \right)
\end{equation}
The upper bound of $10$ on $r$ ensures numerical stability during solver initialization when $\tilde{S}$ may momentarily approach zero. 

It should be explicitly noted that the original formulation by Spalart and Allmaras \cite{Spalart1994} includes additional geometric trip functions ($f_{t1}$ and $f_{t2}$) designed to manually trigger the transition from laminar to turbulent flow at specific user-defined coordinates. However, following the topology optimization methodology established by Yoon \cite{Yoon2016}, these transition functions are deliberately omitted from this thesis. Because the topological boundaries are unknown *a priori* and evolve during optimization, tracking discrete trip points is mathematically unfeasible; therefore, the model assumes a fully turbulent flow regime.

To close the system, the standard empirical calibration constants formulated by Spalart and Allmaras are utilized throughout the FEniCS implementation \cite{NasaTurb}:
\begin{align*}
	c_{b1} &= 0.1355 & \sigma &= 2/3 & c_{b2} &= 0.622 \\
	\kappa &= 0.4187 & c_{w2} &= 0.3 & c_{w3} &= 2.0 \\
	c_{v1} &= 7.1 & c_{w1} &= \frac{c_{b1}}{\kappa^2} + \frac{1 + c_{b2}}{\sigma}
\end{align*}

\subsection{Wall-Distance Calculation via the Relaxed Eikonal Equation}
A major computational challenge in the SA model is the calculation of the wall distance $y$. In a standard static CFD simulation, $y$ is purely geometric and can be pre-computed. However, to ensure this verification solver translates seamlessly into the topology optimization code (where solid boundaries evolve continuously), the implementation abandons discrete geometric searches in favor of a PDE-based approach. 

The true geometric wall distance satisfies the standard Eikonal equation, $|\nabla y| = 1$. However, this formulation contains singularities at the medial axes of the domain. To guarantee continuous differentiability—a strict requirement for later adjoint derivations—a relaxed Eikonal equation proposed by Yoon \cite{Yoon2016} is solved:
\begin{equation}
	\nabla G \cdot \nabla G + \sigma_w G (\nabla \cdot \nabla G) = (1 + 2\sigma_w) G^4
\end{equation}
Here, $G$ is an intermediate, non-dimensional field variable. The parameter $\sigma_w$ is a relaxation constant (typically set to $0.1$) that controls the smoothness of the field. The PDE is subjected to a Dirichlet boundary condition at the physical walls, $G = G_0$, where $G_0$ is an arbitrarily large constant (e.g., $G_0 = 20.0$). Once the field $G$ is solved over the domain, the physical wall distance $y$ is recovered algebraically at every node using the inverse relationship:
\begin{equation}
	y = \frac{1}{G} - \frac{1}{G_0}
\end{equation}

This system is implemented in the core solver suite within the \texttt{Utilities.py} script. As shown in Listing~\ref{lst:eikonal_fenics}, once the non-linear variational problem for $G$ is solved, the physical distance $y$ is recovered and strictly bounded from below to prevent numerical singularities.

\begin{lstlisting}[language=Python, caption={Implementation of Yoon's relaxed Eikonal equation for wall-distance calculation, extracted from \texttt{Utilities.py}.}, label={lst:eikonal_fenics}]
	# Weak form corresponding to Yoon 2016 Eq. (19):
	F = (
	(Constant(1.0) - sigma_w_const) * inner(grad(G), grad(G)) * z * dx
	- sigma_w_const * G * inner(grad(G), grad(z)) * dx
	- (Constant(1.0) + Constant(2.0) * sigma_w_const) * G**4 * z * dx
	)
	problem = NonlinearVariationalProblem(F, G, bcs=wall_bcs, J=derivative(F, G))
	solver = NonlinearVariationalSolver(problem)
	solver.solve()
	
	# Bound G to prevent division by zero
	G.vector()[:] = np.maximum(G.vector()[:], float(g_floor))
	
	# Recover physical wall distance (y)
	y = project(Constant(1.0) / G - Constant(1.0) / g0_const, Space)
	return bound_from_bellow(y, 0.0)
\end{lstlisting}

\section{Variational Formulation and FEniCS Implementation}
The standard Galerkin weak formulation is constructed by multiplying the steady transport Equation~\ref{eq:sa_steady_transport} by a test function $\xi$ and applying integration by parts to the diffusive terms. Because the equation is highly non-linear, robust convergence is heavily dependent on the quality of the mesh and the stability of the external velocity field provided by the Picard outer loop.

The core SA physics, boundary conditions, continuous Eikonal formulations, and the overarching Picard iteration scheme were derived independently for this thesis. However, the underlying FEniCS solver architecture was structurally inspired by an open-source $k-\epsilon$ repository developed by Marcibál \cite{Marcibal2024, joove_keps}, specifically regarding the procedural execution of the variational forms and the implementation of the Incremental Pressure Correction Scheme (IPCS). The explicit translation of the SA mathematical formulation is handled by the \texttt{SpalartAllmarasSteady} class within the \texttt{TurbulenceModel\_SpalartAllmaras.py} script. As demonstrated in Listing~\ref{lst:sa_fenics}, the Unified Form Language (UFL) mirrors the analytical weak form precisely, utilizing the most recent convective velocity field (\texttt{external\_velocity}) passed from the IPCS flow solver.

\begin{lstlisting}[language=Python, caption={FEniCS implementation of the steady-state SA standard Galerkin weak formulation. Python dictionary lookups have been simplified for readability.}, label={lst:sa_fenics}]
	def construct_forms(self, external_velocity):
	self._construct_turbulent_quantities(external_velocity)
	sigma = Constant(2.0 / 3.0)
	
	# Steady-state SA transport equation weak form
	FNT = (
	dot(external_velocity, nabla_grad(self._nu_tilde)) * self._xi * self._dx
	+ inner((self._nu_laminar + self._nu_tilde0) / sigma * grad(self._nu_tilde), grad(self._xi)) * self._dx
	+ self._react_nt * self._nu_tilde * self._xi * self._dx
	- self._source_nt * self._xi * self._dx
	)
	
	self._a_nt = lhs(FNT)
	self._l_nt = rhs(FNT)
\end{lstlisting}

\section{General Solver Configuration and Numerical Parameters}
Before detailing the specific verification benchmarks, it is necessary to establish the general numerical configuration of the FEniCS solver. Both the channel flow and the U-bend verifications were simulated using the segregated algorithmic structure detailed earlier: an outer Picard fixed-point loop couples the fluid momentum equations to the Spalart-Allmaras transport equation, while an inner IPCS resolves the velocity-pressure subsystem via pseudo-time stepping. 

Because the U-bend domain introduces severe geometric curvature, strong convective transport, and adverse pressure gradients, it represents a significantly more demanding numerical challenge than a fully developed, unidirectional channel flow. Consequently, the U-bend simulations necessitated stricter numerical controls to maintain stability, including a heavily refined IPCS pseudo-time step ($\Delta t = 1.0 \times 10^{-4}$ s) and much tighter under-relaxation of the turbulence field updates ($\alpha_{\tilde{\nu}} = 0.01$). To ensure strict scientific reproducibility without disrupting the narrative flow of this verification, the comprehensive matrices detailing all physical boundary conditions, geometric dimensions, and exact numerical solver parameters for both benchmarks have been compiled and are documented extensively in Sections~\ref{sec:app_channel} and \ref{sec:app_ubend} of Appendix~\ref{app:baseline_configurations}.

\section{Verification Benchmark 1: Channel Flow}
To verify the fundamental correctness of the custom fluid and turbulence solvers, the code was first tested on a standard 2D turbulent channel flow. To cleanly isolate and verify turbulence production, wall damping, and Picard coupling without the artifacts of a developing boundary layer, this benchmark was strictly formulated as a pressure-driven periodic channel rather than a conventional inlet-outlet duct. 

The velocity field ($\mathbf{u}$) and Spalart-Allmaras working variable ($\tilde{\nu}$) were constrained periodically between the streamwise boundaries, while the top and bottom boundaries were treated as standard no-slip walls ($\mathbf{u} = 0, \tilde{\nu} = 0$). Flow was driven entirely by a fixed pressure drop from the upstream plane ($p=2$ Pa) to the downstream plane ($p=0$ Pa). Based on a reference bulk velocity of $U_{\mathrm{ref}}=18.5$ m/s and utilizing the full channel height ($H=2.0$ m) as the characteristic length scale, the solver was targeted to resolve a fully turbulent flow regime at a channel Reynolds number of $Re_H \approx 20,300$. By explicitly defining the Reynolds number based on the channel height rather than the mathematical hydraulic diameter, this formulation directly aligns with standard academic conventions for fundamental 2D channel flow verification. To rigorously verify the mathematical accuracy of the formulation, the FEniCS results were compared directly against the commercial Ansys Fluent solver operating under identical periodic boundary conditions.

\subsection{Mesh Resolution Requirements}
Because the SA formulation is a "low-Reynolds" model that does not utilize empirical wall functions, it requires mathematical integration directly through the viscous sublayer. Consequently, standard uniform computational grids are physically insufficient. To satisfy this strict boundary layer requirement, a custom structured inflation mesh was generated for the channel flow. Because the domain is a simple periodic rectangle, this highly anisotropic grid was constructed programmatically utilizing NumPy arrays and converted into a FEniCS-compatible format via the open-source \texttt{meshio} library. Based on the target bulk Reynolds number, this approach explicitly enforced a physical first-cell height of approximately $1.7 \times 10^{-3}$ meters, guaranteeing a non-dimensional wall distance of $y^+ \approx 1$ across the entire domain. 

Furthermore, to guarantee that the numerical results were not influenced by spatial discretization, a formal grid independence study was conducted. The full results of this spatial verification, proving convergence, are documented in Appendix~\ref{subsec:app_channel_mesh}.

\subsection{FEniCS vs. Ansys Comparison}
The channel flow was simulated in FEniCS using the decoupled Picard iteration method and compared against the Ansys Fluent baseline. To validate the solver accuracy both qualitatively and quantitatively, velocity contours were compared, and flow data was extracted along a wall-normal line probe spanning the width of the fully developed channel.

\begin{figure}[htpb]
	\centering
	\begin{subfigure}[b]{0.48\textwidth}
		\centering
		\includegraphics[width=\textwidth]{example-image} % Replace with Ansys channel contour
		\caption{Ansys Fluent SA Contour}
		\label{fig:channel_ansys}
	\end{subfigure}
	\hfill
	\begin{subfigure}[b]{0.48\textwidth}
		\centering
		\includegraphics[width=\textwidth]{example-image} % Replace with FEniCS channel contour
		\caption{FEniCS SA Contour}
		\label{fig:channel_fenics}
	\end{subfigure}
	
	\vspace{0.5cm}
	
	\begin{subfigure}[b]{0.8\textwidth}
		\centering
		\includegraphics[width=\textwidth]{example-image} % Replace with channel velocity profile plot
		\caption{Velocity Profiles}
		\label{fig:channel_profiles}
	\end{subfigure}
	\caption{Verification of the 2D Channel Flow ($Re_H \approx 20,300$). Comparison of the qualitative velocity contours between the Ansys Fluent baseline (a) and the custom FEniCS implementation (b), alongside the quantitative velocity profiles extracted across the channel (c).}
	\label{fig:channel_benchmark}
\end{figure}

As shown in Figure~\ref{fig:channel_benchmark}, the qualitative velocity contours of the custom FEniCS implementation visually replicate the Ansys Fluent baseline. More importantly, the quantitative line plots confirm that the velocity profiles generated by the custom FEniCS model match the Ansys reference data impeccably. The custom solver successfully resolves the logarithmic region of the turbulent boundary layer and accurately predicts the turbulent diffusion toward the channel center. This perfect alignment provides strict mathematical proof that the core Navier-Stokes solver, the relaxed Eikonal field, and the highly non-linear SA production/destruction terms are functioning correctly before introducing the geometric complexities of the U-bend.

\section{Verification Benchmark 2: U-Bend Geometry}
Following the successful channel flow verification, the solver was tested on the highly convective, separating secondary flows of the U-bend geometry evaluated in Chapter~\ref{ch:turbulence_modeling}. It is important to distinguish the physical scope of this specific verification test from the standard Ansys Verification Manual (VMFL048) case. 

While the original VMFL048 case is a full 3D simulation, the verification performed here utilized a strictly modified computational domain to ensure efficiency. First, the domain was reduced to a 2D geometry corresponding directly to the symmetry mid-plane of the 3D case. Second, to reduce unnecessary computational overhead during eventual topology optimization, the highly extended straight inlet and outlet sections ($1.555$ m) of the standard benchmark were explicitly truncated to strictly $30D$ ($0.840$ m).

To drive the flow, a uniform prescribed velocity ($\mathbf{u}_{in} = 1.42$ m/s) was imposed at the inlet alongside a fully turbulent boundary value ($\tilde{\nu}_{in} = 9.34 \times 10^{-5}$ m$^2$/s), while a zero relative static pressure ($p=0$ Pa) was assigned to the outlet. While the FEniCS domain is mathematically a 2D planar approximation of the symmetry mid-plane, the 3D physical pipe diameter ($D = 0.028$ m) and geometric radius of curvature ($R_c = 0.125$ m) were strictly retained as the characteristic length scales. Conserving these parameters ensures that the computed Reynolds number ($Re_D \approx 44,700$) and Dean number ($De \approx 15,000$) identically match the 3D Ansys reference case, maintaining non-dimensional consistency and correctly governing the internal secondary flows.

\subsection{Mesh Resolution Requirements}
Because this formulation utilizes a standard Galerkin weak form without spatial upwind stabilization, numerical stability in highly convective flows is strictly dependent on adequate grid resolution. As established in the channel flow benchmark, the SA model requires integration entirely through the viscous sublayer. 

Unlike the simple rectangular channel where a programmatic NumPy script was sufficient, the U-bend geometry features curved inner and outer arcs, complex corner transitions, and anisotropic boundary-layer inflation around non-orthogonal walls. To manage this severe geometric complexity while satisfying the $y^+ \approx 1$ resolution requirement, a custom mesh was generated using the advanced open-source mesh generator \textbf{Gmsh} \cite{Geuzaine2009}. The first element height adjacent to the wall was strictly controlled to $1.20 \times 10^{-5}$ meters across the entire domain. As documented comprehensively in Appendices~\ref{subsec:app_ubend_truncation} and \ref{subsec:app_ubend_mesh}, strict independence studies were conducted to explicitly prove that neither the $30D$ domain truncation nor the spatial grid discretization artificially influenced the secondary flow physics.

\subsection{FEniCS vs. Ansys Comparison}
Having mathematically validated the spatial resolution and boundary placement, an independent Ansys Fluent simulation was executed on the exact shortened 2D domain. This served as a strict 1:1 baseline to verify the accuracy of the custom FEniCS implementation. 

\begin{figure}[htpb]
	\centering
	\begin{subfigure}[b]{0.48\textwidth}
		\centering
		\includegraphics[width=\textwidth]{example-image} % Replace with Ansys U-bend contour
		\caption{Ansys Fluent SA Contour}
		\label{fig:ubend_ansys}
	\end{subfigure}
	\hfill
	\begin{subfigure}[b]{0.48\textwidth}
		\centering
		\includegraphics[width=\textwidth]{example-image} % Replace with FEniCS U-bend contour
		\caption{FEniCS SA Contour}
		\label{fig:ubend_fenics}
	\end{subfigure}
	
	\vspace{0.5cm}
	
	\begin{subfigure}[b]{0.8\textwidth}
		\centering
		\includegraphics[width=\textwidth]{example-image} % Replace with U-bend velocity profile plot
		\caption{Velocity Profiles}
		\label{fig:ubend_profiles}
	\end{subfigure}
	\caption{Verification of the 2D U-Bend Flow ($Re_D \approx 44,700$, $De \approx 15,000$). Comparison of the qualitative velocity contours between the Ansys Fluent baseline (a) and the custom FEniCS implementation (b), showing successful capture of the separation bubble, alongside the quantitative velocity profiles extracted at the bend exit (c).}
	\label{fig:ubend_verification}
\end{figure}

As shown in Figure~\ref{fig:ubend_verification}, the global flow fields and secondary flow separation bubbles observed in the qualitative velocity contours demonstrate excellent visual agreement. However, Research Question 2 strictly demands \textit{quantitative} validation. To satisfy this requirement, velocity data was extracted along a cross-sectional line probe positioned directly at the bend exit, marking the exact transition from the end of the curvature to the start of the straight outlet section. 

The quantitative velocity profile computed by the custom FEniCS implementation traces the commercial Ansys results closely. The FEniCS implementation accurately replicates the boundary layer growth, the momentum shift toward the outer wall, and the turbulent diffusion predicted by the commercial solver, thereby confirming the mathematical integrity of the code.

\section{Conclusion}
This chapter has detailed the mathematical formulation, stabilization strategy, and rigorous verification of the steady-state Spalart-Allmaras model within the FEniCS framework. By successfully replicating the benchmark results of Ansys Fluent, this implementation directly answers the second research question of this thesis: achieving numerical stability and quantitative agreement in a custom finite element solver, laying a fully consistent foundation for the topology optimization routines developed in the subsequent chapters. 

First, numerical stability was successfully achieved without the need for complex mathematical stabilization artifacts. By decoupling the Navier-Stokes and Spalart-Allmaras equations via a rigorous Picard iteration scheme, and utilizing an Incremental Pressure Correction Scheme (IPCS) to manage the convective fluid terms, the highly non-linear system was robustly driven to equilibrium. Supported by high-density mesh resolution and Yoon's relaxed Eikonal equation, this standard Galerkin formulation completely avoided the numerical oscillations that typically plague turbulent finite-element solvers. 

Second, quantitative agreement was thoroughly verified. The solver produced velocity and eddy viscosity profiles that perfectly matched the commercial Ansys Fluent baseline in both simple periodic channel flows ($Re_H \approx 20,300$) and the highly convective, geometrically complex U-bend ($Re_D \approx 44,700$, $De \approx 15,000$). Consequently, this decoupled Picard FEniCS solver is proven to be a mathematically sound and reliable fluid dynamics engine, specifically structured to support the continuous adjoint topology optimization campaigns presented in Chapters~\ref{ch:topology_optimization_methodology} and \ref{ch:benchmark_results}.